% topic: algorithms
\question กำหนดเหรียญ 3 ชนิด มีมูลค่า \textbf{1}, \textbf{8} และ \textbf{13} บาท โดยมีจำนวนเหรียญแต่ละชนิดไม่จำกัด ต้องการทอนเงินจำนวนหนึ่งให้ได้จำนวนเหรียญ\textbf{น้อยที่สุด}
\newline
\textbf{ตัวอย่าง:} ต้องการทอนเงิน 20 บาท สามารถใช้เหรียญ 13, 1, 1, 1, 1, 1, 1, 1 (รวม 8 เหรียญ) หรือใช้เหรียญ 8, 8, 1, 1, 1, 1 (รวม 6 เหรียญ) หรือใช้เหรียญ 1 ทั้งหมด (20 เหรียญ) เป็นต้น ซึ่งจะเห็นว่าการเลือกเหรียญที่แตกต่างกันให้จำนวนเหรียญรวมที่ไม่เท่ากัน
\newline
\textbf{คำถาม:} ถ้าต้องการทอนเงินจำนวน \textbf{77} บาท ด้วยเหรียญ \{1, 8, 13\} บาท จะต้องใช้เหรียญ\textbf{น้อยที่สุด}กี่เหรียญ
\newline\newline
ตอบ ในกระดาษคำตอบ \newline
% answer: 8
% solution:
%   Greedy (เลือกเหรียญใหญ่สุดก่อน — ไม่ถูกต้อง):
%     77 → 13+13+13+13+13 = 65, เหลือ 12
%        → 8 = 8, เหลือ 4
%        → 1+1+1+1 = 4
%      รวม: 5 + 1 + 4 = 10 เหรียญ ✗
%
%   Optimal (ถูกต้อง):
%     77 → 13×4 + 8×3 + 1×1 = 52 + 24 + 1 = 77
%      รวม: 4 + 3 + 1 = 8 เหรียญ ✓
%
%   ตรวจสอบความเป็น optimal โดยแจกแจงตามจำนวนเหรียญ 13:
%     #13=0: 77 ด้วย {1,8} → 8×9=72 + 1×5 = 14 เหรียญ
%     #13=1: 77-13=64 ด้วย {1,8} → 8×8=64 = 1+8=9 เหรียญ
%     #13=2: 77-26=51 ด้วย {1,8} → 8×6=48 + 1×3 = 2+6+3=11 เหรียญ
%     #13=3: 77-39=38 ด้วย {1,8} → 8×4=32 + 1×6 = 3+4+6=13 เหรียญ
%     #13=4: 77-52=25 ด้วย {1,8} → 8×3=24 + 1×1 = 4+3+1=8 เหรียญ ✓
%     #13=5: 77-65=12 ด้วย {1,8} → 8+1×4 = 5+1+4=10 เหรียญ
%     #13=6: 77-78 < 0 → เป็นไปไม่ได้
%     ดังนั้น จำนวนเหรียญที่น้อยที่สุดคือ 8 เหรียญ (13×4 + 8×3 + 1×1)
